\SourcePage{19}{24}
\section{Photons}

We now turn to a discussion of the field-quantization formulas obtained above.

First of all, formula~(2.12) for the field energy presents the following difficulty. The lowest energy level \SourcePage{20}{25} of the field is attained when the quantum numbers $N_{\mathbf{k}\alpha}$ of all the oscillators vanish (this state is called the \emph{vacuum state of the electromagnetic field}). Yet even in this state every oscillator has a nonzero ``zero-point energy'' $\omega/2$. Summation over the entire infinite set of oscillators therefore gives an infinite result. We thus encounter one of the ``divergences'' arising from the fact that the present theory is not logically complete.

As long as only the eigenvalues of the field energy are at issue, this difficulty can be removed simply by omitting the zero-point energy, that is, by writing the field energy and momentum as\footnote{This omission can be carried out formally without inconsistency if products of operators in~(2.10) are understood to be ``normal'' products, i.e. products in which the operators $\widehat c^+$ always stand to the left of the operators $\widehat c$. Formula~(2.23) then takes the form
\[
\widehat H=\sum_{\mathbf{k}\alpha}\bigl(\omega\widehat c^+_{\mathbf{k}\alpha}\widehat c_{\mathbf{k}\alpha}\bigr).
\]}:
\begin{equation}
E=\sum_{\mathbf{k}\alpha}N_{\mathbf{k}\alpha}\omega,
\qquad
\mathbf{P}=\sum_{\mathbf{k}\alpha}N_{\mathbf{k}\alpha}\mathbf{k}.
\tag{3.1}
\end{equation}

These formulas allow us to introduce a concept fundamental to all of quantum electrodynamics: \emph{light quanta}, or \emph{photons}\footnote{The photon concept was first introduced by Einstein (A.~Einstein, 1905).}. More precisely, the free electromagnetic field may be regarded as an assembly of particles, each having energy $\omega(=\hbar\omega)$ and momentum $\mathbf{k}(=\mathbf{n}\hbar\omega/c)$. The relation between a photon's energy and momentum is just the one required by relativistic mechanics for particles of zero rest mass moving at the speed of light. The occupation numbers $N_{\mathbf{k}\alpha}$ now have the meaning of the numbers of photons with specified momenta $\mathbf{k}$ and polarizations $\mathbf{e}^{(\alpha)}$. A photon's polarization property is analogous to the spin concept for other particles (the special features of the photon in this respect will be considered below, in~\S~6).

It is readily seen that the mathematical formalism developed in the preceding section fully agrees with the interpretation of an electromagnetic field as an assembly of photons; it is precisely the machinery of so-called second quantization applied to a photon system\footnote{The second-quantization method was developed for radiation theory by Dirac (P.~A.~M.~Dirac, 1927).}. In this method (see III, \S~64), the occupation numbers of the states serve as the independent variables, and operators act on functions of these numbers. The principal part is played by particle ``anni\SourcePageJoin{21}{26}hilation'' and ``creation'' operators, which respectively decrease or increase an occupation number by one. The operators $\widehat c_{\mathbf{k}\alpha}$ and $\widehat c^+_{\mathbf{k}\alpha}$ are exactly of this kind: $\widehat c_{\mathbf{k}\alpha}$ annihilates a photon in state $\mathbf{k}\alpha$, while $\widehat c^+_{\mathbf{k}\alpha}$ creates a photon in that state.

The commutation rule~(2.16) is the rule for particles obeying Bose statistics. Hence photons are bosons, as could have been anticipated: any number of photons is allowed in a given state (we shall return in~\S~5 to the significance of this fact).

The plane waves $\mathbf{A}_{\mathbf{k}\alpha}$~(2.26), which occur in the operator $\widehat{\mathbf A}$~(2.17) as coefficients of the photon annihilation operators, can be interpreted as wave functions of photons with definite momenta $\mathbf{k}$ and polarizations $\mathbf{e}^{(\alpha)}$. This interpretation corresponds to expanding the $\psi$-operator in a series of wave functions for the stationary states of a particle in the nonrelativistic second-quantization formalism (unlike that formalism, however, expansion~(2.17) contains both particle annihilation and particle creation operators; the meaning of this difference will become clear later; see~\S~12).

The wave function~(2.26) is normalized by the condition
\begin{equation}
\int \frac{1}{4\pi}\bigl(|\mathbf{E}_{\mathbf{k}\alpha}|^2+|\mathbf{H}_{\mathbf{k}\alpha}|^2\bigr)d^3x=\omega.
\tag{3.2}
\end{equation}
This is normalization to one photon in the volume $V=1$. Indeed, the integral on the left-hand side is the quantum-mechanical expectation value of the photon energy in the state described by the given wave function\footnote{Notice that the factor $1/(4\pi)$ in the integral~(3.2) is twice the usual factor $1/(8\pi)$ in~(2.10). Ultimately, this difference is connected with the fact that the vectors $\mathbf{E}_{\mathbf{k}\alpha}$ and $\mathbf{H}_{\mathbf{k}\alpha}$ are complex, whereas the field operators $\widehat{\mathbf E}$ and $\widehat{\mathbf H}$ are Hermitian.}. The right-hand side of~(3.2), in turn, is the energy of one photon.

Maxwell's equations serve as the photon's ``Schrödinger equation.'' In the present case (for a potential $\mathbf{A}(\mathbf{r},t)$ satisfying condition~(2.1)), this is the wave equation
\[
\frac{\partial^2\mathbf A}{\partial t^2}-\Delta\mathbf A=0.
\]
In the general case of arbitrary stationary states, the photon ``wave functions'' are complex solutions of this equation whose time dependence is supplied by the factor $e^{-i\omega t}$.

In speaking of a photon wave function, we emphasize once more that it must by no means be regarded as the probability amplitude for a photon's \SourcePage{22}{27} spatial localization---in contrast to the basic meaning of a wave function in nonrelativistic quantum mechanics. The reason is that (as pointed out in~\S~1) the very concept of photon coordinates has no physical meaning. We shall return to the mathematical aspect of this circumstance at the end of the next section.

The components of the spatial Fourier expansion of $\mathbf{A}(\mathbf{r},t)$ form the photon wave function in the momentum representation; write it as $\mathbf{A}(\mathbf{k},t)=\mathbf{A}(\mathbf{k})e^{-i\omega t}$. Thus, for a state of definite momentum $\mathbf{k}$ and polarization $\mathbf{e}^{(\alpha)}$, the momentum-representation wave function is simply the coefficient of the exponential factor in~(2.26):
\begin{equation}
\mathbf{A}_{\mathbf{k}\alpha}(\mathbf{k}',\alpha')
=\sqrt{4\pi}\,\frac{\mathbf{e}^{(\alpha)}}{\sqrt{2\omega}}\,
\delta_{\mathbf{k}'\mathbf{k}}\delta_{\alpha\alpha'}.
\tag{3.3}
\end{equation}

Because the momentum of a free particle is measurable, the momentum-representation wave function has a deeper physical meaning than its coordinate-representation counterpart: it permits calculation of the probabilities $w_{\mathbf{k}\alpha}$ of the various possible values of momentum and polarization for a photon in the specified state. By the general rules of quantum mechanics, $w_{\mathbf{k}\alpha}$ is given by the squared modulus of the coefficients in the expansion of $\mathbf{A}(\mathbf{k}')$ in wave functions of states with definite $\mathbf{k}$ and $\mathbf{e}^{(\alpha)}$:
\[
w_{\mathbf{k}\alpha}\propto
\left|\sum_{\mathbf{k}'\alpha'}
\mathbf{A}^*_{\mathbf{k}\alpha}(\mathbf{k}',\alpha')\mathbf{A}(\mathbf{k}')
\right|^2
\]
(the proportionality factor depends on the normalization convention for the functions). Substitution of~(3.3) gives
\begin{equation}
w_{\mathbf{k}\alpha}\propto
\left|\mathbf{e}^{(\alpha)}\mathbf{A}(\mathbf{k})\right|^2.
\tag{3.4}
\end{equation}
After summing over both polarizations, we obtain the probability that the photon has momentum $\mathbf{k}$:
\begin{equation}
w_{\mathbf{k}}\propto|\mathbf{A}(\mathbf{k})|^2.
\tag{3.5}
\end{equation}
